\documentclass[10pt,a4paper]{amsart}
\usepackage[margin=19mm]{geometry}
\usepackage{amsmath,amssymb,mathtools}
\usepackage[scr=rsfs,cal=euler]{mathalfa}
\usepackage{xfrac}
\usepackage[unicode,hidelinks,bookmarks=false]{hyperref}
\newcommand{\ie}{i.e.\ }
\newcommand{\cf}{cf.\ }
\newcommand{\resp}{respectively\ }
\newcommand{\Subsection}{\subsection}
\providecommand{\nobreakdash}{\nobreak\hbox{-}}
\setlength{\emergencystretch}{2em}
\multlinegap0pt
\usepackage{amssymb}
\usepackage{bm}
\usepackage{mathtools}
\newtheorem{proposition}{Proposition}
\newtheorem{lemma}{Lemma}
\def\Hs{\mathbf H}
\def\Is{\mathbf I}
\def\NG{{\tt G}}
\def\Ps{\mathbf P}
\def\Rs{\mathbf R}
\def\S{{\mathbf S}}
\title{A generalized cluster structure on $GL_n$ via birational Poisson maps}
\author{Misha Gekhtman, Michael Shapiro, Alek Vainshtein}
\date{Source: arXiv:2603.12486v1 (CC BY 4.0)}
\begin{document}
\maketitle

\noindent\emph{Source and licence.} A generalized cluster structure on $GL_n$ via birational Poisson maps. Version arXiv:2603.12486v1 (CC BY 4.0), distributed by arXiv under the Creative Commons Attribution 4.0 International licence, http://creativecommons.org/licenses/by/4.0/, as recorded in the arXiv licence metadata for this exact version and shown on the abstract page https://arxiv.org/abs/2603.12486v1. Full licence notice: \texttt{licenses/arxiv-2603.12486-CC-BY-4.0.txt}.\par\smallskip

\noindent\textit{Editorial scope.} Counted unit: the article's lemma lemma (source0.tex lines 3138--3140). The article's complete original proof of it is delivered with it (source0.tex lines 3142--3245, one \texttt{proof} environment of 104 lines). No mathematics is written, repaired or re-derived: the statement and the proof are the article's own lines, in the article's own notation, and no retained line is substituted or re-flowed.  Retained uncounted as the source-local context the proof needs: lines 3049--3051; lines 3055--3059; lines 3062--3067; lines 3079--3082; lines 3086--3090.  One declared adaptation: exactly 3 ancillary strings inside the retained spans are omitted (image-inclusion commands and, where the source repeats a label, the repeated label definitions) and no other line differs from the source; the figure environments, with their placement, captions and labels, are kept, and the package records the omitted strings in fidelity receipts bound to the affected bodies.  No retained line contains a citation, so the wrapper carries no bibliography and no bibliography entry.  The retained lines contain the article's own diagram code, which the wrapper typesets with the standard tikz packages; no image file is delivered and the package declares no asset dependency.  Editorial formatting only, in addition: an AMS \texttt{amsart} preamble that restates the article's own declarations and macros used by the retained lines, namely the environments \texttt{proposition}, \texttt{lemma} on separate counters, together with the standard packages those lines need.  The retained environments print in this document as Proposition 1, Lemma 1 in order of appearance, whereas the article prints the same items with the numbering of its own full document.  The complete native source of the article as distributed in the arXiv e-print for this exact version is bundled unchanged as \texttt{sources/196280/source0.tex}.
\begin{equation}\label{4pluck}
[k-1\,2]\cdot[\overline{k-1}]=[\overline{k-2}\,2]\cdot[k]+[\overline{k-1}\,2]\cdot[k-1].
\end{equation}

\begin{proposition}\label{promo1}
{\rm (i)} The four-term Pl\"ucker relation for the promotion of a matrix $\NG(k_1,\dots,k_{p-1},2)$ is obtained by promoting every function involved in the four-term Pl\"ucker relation for $\NG(k_1,\dots,k_{p-1},2)$.

{\rm(ii)} The four-term Pl\"ucker relation for the demotion of a promoted matrix is obtained by demoting every function involved in the four-term Pl\"ucker relation for the promoted matrix.
\end{proposition}

\begin{equation}\label{promo4pluck}
\begin{aligned}
&[2^q\,k-1\,2]\cdot[\bar1\, 2^{q-1}\,k]=[\bar1\,2^{q-1}\,k-1\,2]\cdot[2^q\,k]+[\bar1\,2^{q-1}\,k\,2]\cdot[2^q\,k-1],\\
&[\bar1\,2^{q-1}\,k-1\,2]\cdot[2^{q-1}\,k]=[2^{q-1}\,k-1\,2]\cdot[\bar1\,2^{q-1}\,k]+[2^{q-1}\,k\,2]\cdot[\bar1\,2^{q-1}\,k-1].
\end{aligned}
\end{equation}

\begin{equation}\label{DJ}
[k_1-1\,k_2]\cdot[\overline{k_1-1}\,k_2-1]=[k_1\,k_2-1]\cdot[\overline{k_1-2}\,k_2]+[\overline{k_1-1}\,k_2]
\cdot[k_1-1\,k_2-1].
\end{equation}

\begin{proposition}\label{promo2}
{\rm (i)} The Desnanot--Jacobi identity for the promotion of a matrix $\NG(k_1,\dots,k_{p})$ is obtained by promoting every function involved in the Desnanot--Jacobi identity for $\NG(k_1,\dots,k_p)$.

{\rm (ii)} The Desnanot--Jacobi identity for the demotion of a promoted matrix is obtained by demoting every function involved in the Desnanot--Jacobi identity for the promoted matrix.
\end{proposition}

\begin{lemma} Mutation sequence $\Hs_n$ applied to the quiver $Q_n^0$ results in $Q_n^n$. The variable attached to vertex $|i,j|$ in $Q_n^n$ is $[\overline{j-2}\ n-i-j+4]$ for $j\ge 3$ except for vertex $C=|1,3|$, $[n-i+2]$ for 
$j=2$, and $[\overline{n-i+2}]$ for $j=1$. %the function attached to $C$ is $x_{n1}$.
\end{lemma}

\begin{proof} The proof is by induction on $n$. For $n=4$, the exchange relations along the mutation path $\Hs_4$ are
\begin{equation}\label{exchange4}
\begin{aligned}
0&:\quad [2^3]\cdot[\bar 1\,3]=[\bar 1\,2^2]\cdot[2\,3]+[2^2]\cdot[\bar 1\,3\,2], \\
-1&: \quad [\bar 1\,2^2]\cdot[3]=[\bar 1\,3]\cdot[2^2]+[3\,2]\cdot[\bar 1\,2],\\
1&: \quad [2\,3]\cdot[\bar 2\,2]=[\bar 1\,3]\cdot[3\,2]+[2^2]\cdot[\bar 2\,3], \\
2&: \quad [3\,2]\cdot[\bar 3]=[\bar 2\,2]\cdot[4]+[3]\cdot[\bar 3\,2],\\
-2&: \quad [2^2]\cdot[\bar 2]=[3]\cdot[\bar 1\,2]+[\bar 2\,2]\cdot[2].
\end{aligned}
\end{equation}
A direct check shows that the assertions of the lemma hold true. Additionaly, one can check that for $n>4$ the above mutation sequence adds the arrow from $3$ to $C$ and deletes the arrow from $5$ to $3$.

Note that the first two exchange relations above are 
particular cases of~\eqref{promo4pluck} for $k=3$ and $q=2$,   
the third relation is~\eqref{DJ} for $k_1=k_2=3$, while the last two are~\eqref{4pluck} for $k=4$ and $k=3$, respectively. We thus see that to handle the case $n=4$ we have to consider matrices $\NG(3,2)$ (more exactly, its promotion $\NG(2,3,2)$), $\NG(3,3)$, and $\NG(4,2)$.

Assume that the assertion of the lemma holds true for $\Hs_{n-1}$, and consider the evolution of $Q_n^0$ and the associated cluster variables along the mutation sequence $\Hs_n$ divided into seven consecutive segments: $\Hs_n=\Ps_{n-1}\Is_{n-1}\Rs_{n-1}\S_{n-1}\Is_n\Rs_n\S_n$. Recall that the promotion property for $Q_{n-1}^0$ and $Q_n^0$ is valid at every mutable vertex of $Q_{n-1}^0$. Moreover, mutations in the first two segments of $\Hs_n$ are identical to mutations in the first two segments of $\Hs_{n-1}=\Ps_{n-1}\Is_{n-1}\Rs_{n-1}\S_{n-1}$, that is, involve vertices with the same universal numbers connected with the same arrows. Therefore, by Propositions~\ref{promo1} and~\ref{promo2},
the promotion property remains valid for all mutated vertices along the sequence.

Further, compare mutations in the third segment of $\Hs_{n-1}$ and the third segment of $\Hs_{n}$. Recall that these
are mutations in the first from the left mutable layer of $Q_{n-1}^0$ and the second from the left mutable layer of $Q_n^0$. The bottom vertex of this layer has universal number $i=(n-3)(n-4)/2$.
The arrows inside this layer and to the right of it are identical in both mutated quivers, as explained above. 
An additional inductive assumption for the quiver obtained from $Q_n^0$ claims the following: 

(i) the arrow from the top vertex to $i$ is deleted; 

(ii) all other arrows inside this layer are the same as in $Q_{n}^0$;

(iii) the arrow from $i$ to $i-1$ is reversed; 

(iv) there are arrows  from $i$ to $C$ and from $C$ to $i-1$. 

\noindent The arrows on the left of this layer are nor involved in the previous mutations, so they remain exactly as in  $Q_{n-1}^0$ and $Q_n^0$, respectively. This assumption can be verified by direct observation for $n=5$.

We start the third segment with the mutation at $i$. The left neighborhood of $i$ in the quiver obtained from $Q_{n-1}^0$ before and after the mutation is shown in Fig.~\ref{ileft}a). The variable attached to the frozen vertex 
$j=(n-2)(n-3)/2+2$ is $[\bar2\,n-2]$. A similar neighborhood in the quiver obtained from $Q_{n}^0$ is shown in 
Fig.~\ref{ileft}b). The variable attached to the mutable vertex $j-1$ is $[3\,n-2]$. Additionally, in the second case, there is $x_{1n}$ attached to vertex $C$. Note that the product of variables attached to $j-1$ and $C$ is $[\bar1\,3\,n-2]$, that is, the promotion of  $[\bar2\,n-2]$. Consequently, Proposition~\ref{promo2} guarantees that the promotion property at $i$ remains valid after this mutation. Moreover, the situation before the mutation at $i+1$ is exactly the same as before, up to shifting all numbers up by one. Consequently, the promotion property at all mutable vertices of 
$Q_{n-1}^0$ remains valid after the third segment is applied. Note that in the last mutation in this segment, the vertices that correspond to $i+1$ and $j-1$ in Fig.~\ref{ileft}b) are the bottom and the top vertices of the next layer, respectively, so that deletion of the arrow from $j-1$ to $i+1$ proves inductive assumption (i), reversal 
of the arrows from $i+1$ to $i$ and from $i$ to $C$ and addition of the arrow from $i+1$ to $C$ proves inductive assumptions (iii) and (iv). All other mutations of this segment do not affect arrows in the leftmost mutable layer of 
$Q_n^0$, which proves inductive assumption (ii). 

\begin{figure}[ht]
\begin{center}

\caption{Left neighborhood of $i$ before and after the first mutation in $\Rs_{n-1}$: a) for $\Hs_{n-1}$; 
b) for $\Hs_n$}
\label{ileft}
\end{center}
\end{figure}

Finally, mutations in the fourth segment in $\Hs_{n-1}$ and $\Hs_n$ again involve only vertices with the same numbers connected with the same arrows, and hence by the end of the first four segments in $\Hs_n$ the promotion property remains valid for all mutable vertices of $Q_{n-1}^0$. Taking into account the induction assumption concerning the 
structure of $Q_{n-1}^{n-1}$, we can describe the quiver $Q_n^{n-1}$ obtained from $Q_n^0$ via the mutation sequence 
$\Hs_{n-1}$ and the variables attached to its vertices, as compared to $Q_{n-1}^{n-1}$. 

Quiver $Q_n^{n-1}$ is placed on the grid, same as $Q_{n-1}^{n-1}$. The restriction of $Q_n^{n-1}$ to positions $|i,j|$, 
$i\in[3,n+1]$, $j\in[1,n-i+2]$ coincides with restriction of $Q_{n-1}^{n-1}$ to positions $|i,j|$, $i\in[2,n]$, $j\in[1,n-i+1]$, so that the universal number for the vertex at position $|i,j|$ in $Q_n^{n-1}$ is the same as for $|i-1,j|$ in $Q_{n-1}^{n-1}$. Consequently, by the promotion property, the variable attached to position $|i,j|$ for $i\in[3,n+1]$, 
$j\in[1,n-i+2]$ is $[\bar1\, j-1\, n-i-j+5]$ for $j\ge 3$, $[2\, n-i+3]$ for $j=2$, and $[\bar1\, n-i+2]$ 
for $j=1$. The only difference between the two restrictions is that  in $Q_{n-1}^{n-1}$ vertices at positions $|2,2|$ and $|n,1|$ are frozen, while in $Q_n^{n-1}$ the corresponding vertices at $|3,2|$ and $|n+1,1|$ are mutable. Frozen vertices from the leftmost layer of $Q_n^0$ are placed at positions $|1,j|$ for $j\in[3,n]$ in the same vertical order. Vertices of the next layer are placed at positions $|2,j|$ for $j\in[2,n]$ in the same vertical order, except for the bottom mutable vertex $(n-2)(n-3)/2$ that has been already placed at $|3,2|$. The arrows between the vertices in these two layers are not affected by the mutations in $\Hs_{n-1}$, except for the cases described in induction assumptions 
(i)--(iv) that are already proved for this layer. Evolution of arrows between this layer and the layer to the right of it can be seen in Fig.~\ref{ileft}b). Finally, vertices $-2n+4$ and $-2n+5$ are placed at positions $|n+1,3|$ and $|n+2,2|$, respectively. Arrows between them and the rest of the vertices are not affected by mutations in $\Hs_{n-1}$. The quiver $Q_6^5$ is shown in Fig.~\ref{Q65}a). The part of $Q_6^5$ that coincides with the described above part of $Q_5^5$ is contained inside the triangle shown by a thin line.


\begin{figure}[ht]
\begin{center}

\caption{Quiver $Q_6^5$: a) the quiver; b) the part of $Q_6^5$ relevant for the sequence $\Is_6$}
\label{Q65}
\end{center}
\end{figure}

By the promotion property, the variables inside this triangle are promotions of the 
corresponding variables in $Q_5^5$, so that by the inductive assumption, the variable attached to $|i,j|$ is
$[\bar1\, j-1\,n-i-j+3]$. The variables outside the triangle retain their initial values. It is easy to see 
that vertices that are mutated in the sequence $\Is_n$ fill the triangle $i\in[3,n-1]$, $j\in[3,n-i+2]$ and that order is by columns from left to right, bottom up within each column. The corresponding vertices for the quiver $Q_6^5$ are shown by double circles in Fig.~\ref{Q65}.  Consequently, the relevant vertices outside the triangle 
are $|2,j|$ with $j\in[3,n-1]$, and the attached variables are $[j\,n-j+1]$. The part of $Q_6^5$ relevant for the sequence $\Is_6$ is shown in Fig.~\ref{Q65}b).

Let us rearrange the vertices of $Q_n^{n-1}$ 
after the sequence $\Is_n$ is applied in the following way: the vertices in rows $j=1,2$ are shifted left by one, and the vertices $|i,j|$ on the upper diagonal $i+j=n+2$ are shifted down by $j$ and right by one, so that they now form the lower row of the quiver. The result of this rearrangement for the quiver $Q_6^5$ is shown in Fig.~\ref{Q65I}.

\begin{figure}[ht]
\begin{center}

\caption{Quiver $Q_6^5$ after mutation sequence $\Is_6$ and rearrangement}
\label{Q65I}
\end{center}
\end{figure}
 
Comparing quivers $Q_n^{n-1}$ (before the application of $\Is_n$) and $Q_{n-1}^{n-2}$ after the application of $\Is_{n-1}$ we observe that the subquiver of the latter quiver that is relevant for the sequence $\Rs_{n-1}\S_{n-1}$ is isomorphic to 
the subquiver of the former one that is relevant for the sequence  $\Is_n$. Moreover, the variable attached in the latter 
quiver to the vertex that correspond to the vertex $|i,j|$ of $Q_n^{n-1}$ is $[j-1\, n-i-j+3]$. It is easy to see that
the variable $[\bar1\, j-1\,n-i-j+3]$ attached to $|i,j|$ in $Q_n^{n-1}$ is obtained from $[j-1\, n-i-j+3]$ by consecutive promotion and demotion, and the same holds true for all other vertices in the two isomorphic subquivers. 
So, by the promotion and demotion properties, we conclude that all exchange relations along the sequence $\Is_n$ are
obtained by consecutive promotion and demotion of the exchange relations along the sequence $\Rs_{n-1}\S_{n-1}$. Therefore, the variable attached to the vertex $|i,j|$ of the rearranged quiver is $[j-1\,n-i-j+4]$ for $i\in[2,n-1]$, 
$j\in [3,n-i+2]$, $[\bar1\, n-i+2]$ for $j=2$, $i\in[3,n]$,  $[n-i+3]$ for $i\in[3,n+1]$, and 
$[\overline{j-2}\, n-j+3]$ for $i=1$, $j\in[4,n+1]$.

Further, the sequence $\Rs_n$ is applied. It is easy to see that the vertices that participate in this sequence are the vertices in the second column of the current quiver bottom up (vertices $6, 7, 8, 9$ in Fig.~\ref{Q65I} for $n=6$). The corresponding exchange relations are~\eqref{DJ} for
$k_1=j$ and $k_2=n-j+2$ with $j\in[3,n-1]$, and~\eqref{4pluck} for $k=n$. The resulting variables are $[\overline{j-1}\,
n-j+1]$ for $j\in[3,n-1]$ and $[\overline{n-1}]$ for $j=n$.

Finally, the sequence $\S_n$ is applied. It is easy to see that the vertices that participate in this sequence fill the triangle $i\in[3,n-1]$, $j\in[3,n-i+2]$ (it is shown by double circles in Fig.~\ref{Q65I}). The order of mutations is 
by columns from left to right, bottom up in each column. Direct inspection shows that the relevant part of the quiver is exactly the same as for sequences $\Rs_{n-1}\S_{n-1}$ and $\Is_n$, see Fig.~\ref{Q65}b) for the case $n=6$. Moreover, 
the functions attached to the vertices involved in the mutations are exactly the same as in the case of $\Rs_{n-1}\S_{n-1}$, and hence the variables attached to the vertices after the mutation coincide with those obtained at the end of $\Hs_{n-1}$. It remains to rearrange the resulting quiver similarly to what was done before: the vertices in the rows $j=1,2$ 
are shifted left by one, the vertices of the first column are shifted down by one,  and the vertices $|i,j|$ on the upper diagonal $i+j=n+2$ are shifted down by $j$ and right by one, so that they now form the lower row of the quiver. The result of this rearrangement coincides with the quiver $Q_n^n$.
\end{proof}

\end{document}
